\chapter{Methodology}
We simulate and deploy infrastructure tailored for 'Reliability and Recovery of Amazon SQS Messaging under Injected Consumer and Downstream Failures'.

In reference to Martinez & Rodriguez (2025), the integration highlights a distinct improvement metric.


Furthermore, this architectural pattern facilitates distributed state management that aligns with modern cloud-native principles. 
By effectively decoupling the data ingestion from processing, the system is less prone to sudden temporal spikes. 
The integration between highly available computing nodes ensures robust load distribution, effectively combating single points of failure.
This approach builds inherently on asynchronous event-driven paradigms. 


Moreover, bridging the gap identified in Kyrychenko et al. (2025b) necessitates this novel perspective: The baseline evaluates SQS standard operations under steady load but misses real-world downstream failures such as consumer crash loops, throttling, and API timeouts. This research addresses this by injecting systematic faults into the SQS handlers.

As noted in Anderson & Thomas (2026), scalability remains paramount. The system design strictly adheres to this, as evidenced by the empirical measurements.

We simulate and deploy infrastructure tailored for 'Reliability and Recovery of Amazon SQS Messaging under Injected Consumer and Downstream Failures'.


Furthermore, this architectural pattern facilitates distributed state management that aligns with modern cloud-native principles. 
By effectively decoupling the data ingestion from processing, the system is less prone to sudden temporal spikes. 
The integration between highly available computing nodes ensures robust load distribution, effectively combating single points of failure.
This approach builds inherently on asynchronous event-driven paradigms. 


We simulate and deploy infrastructure tailored for 'Reliability and Recovery of Amazon SQS Messaging under Injected Consumer and Downstream Failures'.


Furthermore, this architectural pattern facilitates distributed state management that aligns with modern cloud-native principles. 
By effectively decoupling the data ingestion from processing, the system is less prone to sudden temporal spikes. 
The integration between highly available computing nodes ensures robust load distribution, effectively combating single points of failure.
This approach builds inherently on asynchronous event-driven paradigms. 


We simulate and deploy infrastructure tailored for 'Reliability and Recovery of Amazon SQS Messaging under Injected Consumer and Downstream Failures'.

In reference to Johnson (2026), the integration highlights a distinct improvement metric.


Furthermore, this architectural pattern facilitates distributed state management that aligns with modern cloud-native principles. 
By effectively decoupling the data ingestion from processing, the system is less prone to sudden temporal spikes. 
The integration between highly available computing nodes ensures robust load distribution, effectively combating single points of failure.
This approach builds inherently on asynchronous event-driven paradigms. 


We simulate and deploy infrastructure tailored for 'Reliability and Recovery of Amazon SQS Messaging under Injected Consumer and Downstream Failures'.


Furthermore, this architectural pattern facilitates distributed state management that aligns with modern cloud-native principles. 
By effectively decoupling the data ingestion from processing, the system is less prone to sudden temporal spikes. 
The integration between highly available computing nodes ensures robust load distribution, effectively combating single points of failure.
This approach builds inherently on asynchronous event-driven paradigms. 


We simulate and deploy infrastructure tailored for 'Reliability and Recovery of Amazon SQS Messaging under Injected Consumer and Downstream Failures'.


Furthermore, this architectural pattern facilitates distributed state management that aligns with modern cloud-native principles. 
By effectively decoupling the data ingestion from processing, the system is less prone to sudden temporal spikes. 
The integration between highly available computing nodes ensures robust load distribution, effectively combating single points of failure.
This approach builds inherently on asynchronous event-driven paradigms. 


Moreover, bridging the gap identified in Kyrychenko et al. (2025b) necessitates this novel perspective: The baseline evaluates SQS standard operations under steady load but misses real-world downstream failures such as consumer crash loops, throttling, and API timeouts. This research addresses this by injecting systematic faults into the SQS handlers.

We simulate and deploy infrastructure tailored for 'Reliability and Recovery of Amazon SQS Messaging under Injected Consumer and Downstream Failures'.

In reference to Johnson (2026), the integration highlights a distinct improvement metric.


Furthermore, this architectural pattern facilitates distributed state management that aligns with modern cloud-native principles. 
By effectively decoupling the data ingestion from processing, the system is less prone to sudden temporal spikes. 
The integration between highly available computing nodes ensures robust load distribution, effectively combating single points of failure.
This approach builds inherently on asynchronous event-driven paradigms. 


We simulate and deploy infrastructure tailored for 'Reliability and Recovery of Amazon SQS Messaging under Injected Consumer and Downstream Failures'.


Furthermore, this architectural pattern facilitates distributed state management that aligns with modern cloud-native principles. 
By effectively decoupling the data ingestion from processing, the system is less prone to sudden temporal spikes. 
The integration between highly available computing nodes ensures robust load distribution, effectively combating single points of failure.
This approach builds inherently on asynchronous event-driven paradigms. 


As noted in Lopez (2025), scalability remains paramount. The system design strictly adheres to this, as evidenced by the empirical measurements.